% Nouveau EGA V (cote 25): archive cote 25, site page 175 (PDF page 176).
% English translation of the checked French annotated master.
% AI-translated by GPT-6 Astra (OpenAI, ChatGPT Pro mode).
\scanpage{25/175}
\mspage{\hand{2.}}

associated with Q (defined by the matrix of the $\frac{\hand{\partial} Q}{\hand{\partial} X_i}$ ) is alternating, so its
``kernel'' $N\hand{\subset}\overset{\hand{\vee}}{E}$ is such that $\overset{\hand{\vee}}{E}/N$ has \uline{even} rank, hence if E itself
has even rank, the same holds for N. Thus if $N\hand{\neq} 0$, the rank of N
is at least 2, from which it follows that there is at least one nontrivial zero of
Q in N, i.e. Q is not smooth. NB This recourse to coordinates is decidedly
unworthy of us; we should use them only to prove
simply the
\uline{Lemma} \var{\inserted[typed]{16.3.}}{16.3.}{insertion}{ega25_25_175-v0001} E a vector space of finite rank over a field k, \var{\inserted[typed]{$Q\hand{\in} \mathrm{Sym}^2(E)$}}{$Q\hand{\in} \mathrm{Sym}^2(E)$}{insertion}{ega25_25_175-v0002} x a nonzero element of
$\overset{\hand{\vee}}{E}$ , \var{\xout{such that Q(x\unsure{)}}}{}{typewriter strike-out}{ega25_25_175-v0003} x' its image in P(E), suppose $Q(x,x) = 0$, i.e.
$x'\hand{\in} V(Q)$. Then, for \var{\xout{x' to be a smooth point of}}{}{typewriter strike-out}{ega25_25_175-v0004} V(Q) to be smooth at x',
and of \var{\xout{relative}}{}{typewriter strike-out}{ega25_25_175-v0005} dim rank E - 1 (i.e. $Q\hand{\neq} 0$), it is necessary and sufficient that x not
belong to the kernel of the homomorphism $\overset{\hand{\vee}}{E}\hand{\longrightarrow} E$ defined by the symmetric bilinear
form B associated with Q\hand{/}.
16.4. The following study is intended essentially to \var{\xout{treat the case}}{}{typewriter strike-out}{ega25_25_175-v0006}
give a criterion of smoothness for a quadratic form in the case not
covered by 16.2. , i.e. essentially in the case of a vector space of
odd dimension over a field of char 2.\var{\xout{in which}}{}{typewriter strike-out}{ega25_25_175-v0007} \var{\inserted[typed]{In this}}{In this}{insertion}{ega25_25_175-v0008} case \uline{every} quadratic
form Q on $\overset{\hand{\vee}}{E}$ is degenerate, but one readily sees (taking in particular
% NOTE: The canceled typed line stands below "is degenerate" in the source.
\var{\xout{and below}}{}{typewriter strike-out}{ega25_25_175-v0009}
the ``standard form'') that it may very well be smooth.

For every integer $n\hand{>}0$, let us introduce the ``standard quadratic form''
$Q_n$ on $\underline{Z}^n$ , as a form with integer coefficients in n variables
$X_1,..,X_n$ , distinguishing the two cases:
a) $n = 2m$ , $Q_{2m}(X_1,..,X_{2m}) = X_1X_2 + X_3X_4 + \ldots$
b) $n = 2m+1$, $Q_{2m+1}(X_1,..,X_{2m+1}) =$\var{\xout{\unsure{$(X_1^{2}X_2+X$}}}{}{typewriter strike-out}{ega25_25_175-v0010} $Q_{2m}(X_1,..,X_{2m}) + X_{2m+1}^2$.
\uline{Lemma} 16.5 Let n be an odd integer, and consider
\[ \unsure{Q|M}(X_{\unsure{1|n}},..,X_n) = \hand{\sum_i} a_iX_i^{\,2} + \hand{\sum}_{i\hand{<}j} b_{ij}X_iX_j \ , \]
a quadratic form in n indeterminates $X_i$ , with indeterminate coefficients
$a_i$, $b_{ij}$, so that the discriminant d(Q) is a polynomial with

\mspage{\hand{175}}
\tnote{25-175-1}{The source's spelling standart is translated standard. Dim and char mean dimension and characteristic. The displayed Q(x,x), the dimension rank E - 1, and the unfinished canceled formulas remain as written. The canceled relative precedes dimension in English, corresponding to French dimension relative.}
